% chapter/08_extended_analysis.tex
% Session 8 — Extended Analysis
% Compiled with pdflatex + bibtex (nonstopmode)

\section{Extended Analysis}
\label{sec:analysis}

This section extends the core results with ablation studies, oracle analysis, and
efficiency measurements. The goal is to isolate the mechanisms behind the non-monotone
finding and characterise the operating regime where DCA-Trie provides value.

\subsection{Ablation: Gates Are Not the Source of v2's Collapse}
\label{sec:ablation:gates}

Table~\ref{tab:ablation_gates} reports the decisive ablation: DCA-Trie v2 with both
TypeOracle gates disabled.

\begin{table}[htpb]
  \centering
  \caption{Gate ablation on the controlled WebQSP subset.}
  \label{tab:ablation_gates}
  \begin{tabular}{@{}lccc@{}}
    \toprule
    Configuration     & Hits@1 & BUR & max volatility \\
    \midrule
    GCR baseline      & 85.7\% & --- & --- \\
    DCA-Trie v2       & 60.7\% & 0.440 & 0.000 \\
    DCA-Trie v2 (no gates) & 60.7\% & 0.463 & 0.000 \\
    \bottomrule
  \end{tabular}
\end{table}

Disabling both gates leaves the score unchanged at 60.7\%, and only 11 of 300
predictions differ at all (96.3\% byte-identical; McNemar $p = 1.0$). The 25-point gap
versus the baseline is therefore architectural, caused by the per-hop beam-wise trie
rebuilding and the path-length loop, not by the type/range gates. This contradicts the
hypothesis that the DoG-proxy gates are what cost accuracy.

\subsection{Oracle Analysis: Trie Volatility Is Not the Failure Mechanism}
\label{sec:analysis:volatility}

Table~\ref{tab:mechanism} reports the per-question mechanism metrics for v2 and its
no-gates variant.

\begin{table}[htpb]
  \centering
  \caption{Mechanism metrics, v2 versus v2-no-gates (per-question means).}
  \label{tab:mechanism}
  \begin{tabular}{@{}lcc@{}}
    \toprule
    Metric & v2 dynamic & v2 no-gates \\
    \midrule
    Beam utilization ratio (BUR) & 0.440 & 0.463 \\
    BUR entropy                  & 0.738 & 0.823 \\
    SIR decay slope              & $+0.0285$ & 0.0000 \\
    Max constraint volatility    & 0.000 & 0.000 \\
    Max relative volume          & 0.205 & 0.000 \\
    \bottomrule
  \end{tabular}
\end{table}

Beams collapse to roughly 2.2 distinct terminals of 5 in both runs, so diversity
collapse is present but identical with and without the gates. Trie volatility is zero in
both runs: the per-hop tries are stable, and the rebuild-instability hypothesis is
falsified as the failure mechanism. The hop profile is consistent: hop 1 starts with one
beam and expands to five, hop 2 holds five beams, and the trie holds roughly 405 to 455
candidates per hop. The evidence points instead to tokenization misalignment in the
per-hop trie construction as the root cause: entity boundaries that the tokenizer splits
differently in-context than in isolation, so valid tokens are mis-locked at each hop.

\subsection{Efficiency: The Lazy v3 Variant}
\label{sec:efficiency}

Table~\ref{tab:lazy} reports the construction cost of the lazy v3 variant.

\begin{table}[htpb]
  \centering
  \caption{Lazy construction cost per question (controlled WebQSP subset).}
  \label{tab:lazy}
  \begin{tabular}{@{}lcc@{}}
    \toprule
    Metric & Static (baseline/v1) & Lazy (v3) \\
    \midrule
    Candidates materialised & 2,637 (avg) & 754.7 (avg, max 3,480, min 29) \\
    DFS runs & full enumeration & 0 in every question \\
    Frontier builds & --- & 18.4 (avg) \\
    \bottomrule
  \end{tabular}
\end{table}

The static trie enumerates the full DFS path set before decoding; v3 renders only what
the beams reach, materialising an average of 754.7 candidates per question versus 2,637
for the static trie, roughly a 3.5$\times$ reduction in construction cost, with no DFS
performed at all. Construction work therefore tracks the actual decoding trajectory
rather than the worst-case path space.

\subsection{Calibration: The Lazy Mechanism Is Faithful}
\label{sec:calibration}

A sceptical reader could attribute v3's 81.0\% to a leakier constraint rather than a
faithful one. The calibration run answers this. With both TypeOracle gates disabled,
v3 reproduces the GCR baseline exactly: 257/300 (85.7\%) on both conditions, with
100\% per-question agreement. With the gates off, the lazy constraint admits precisely
the same language as the static baseline trie, question for question. The mechanism is
faithful; every gated-number difference is attributable to the gates themselves, not to
the lazy decoding.

\subsection{Summary}
\label{sec:analysis:summary}

The extended analysis isolates three findings behind the non-monotone result.

First, the accuracy cost is attributable to the gates, not the mechanism. Calibration
shows the ungated lazy constraint reproduces the baseline exactly, so the 4.0pp (v1)
and 4.7pp (v3) drops are entirely explained by the type/range gates.

Second, the v2 collapse is architectural. The 25-point gap is unchanged when both gates
are disabled, only 11 of 300 predictions differ, and trie volatility is zero. Step-wise
trie rebuilding itself, not semantic filtering, is the failure mode.

Third, static and lazy type-gating are equivalent in accuracy ($p = 0.77$, ns): the
lazy frontier matches the static trie at roughly 3.5$\times$ lower construction cost,
in a single decoding pass. The lazy mechanism is the efficient way to pay constraint
construction cost only where decoding actually goes.

These findings are not specific to the TypeOracle. Any constraint oracle that removes
paths will exhibit similar trade-offs, and any step-wise trie-rebuilding decoder faces
the architectural risk that v2 demonstrates. The non-monotone result is a structural
property of constrained beam search, not an artifact of our particular oracle design.